%% ===================================================================================
%% Capítulo 7 -- Verificação, Validação e Testes
%% Define a estratégia, os critérios e os casos de teste (CT) rastreados aos requisitos.
%% Os RESULTADOS de execução entram no Capítulo 8. Dono: Arquiteto/Redator; cada agente
%% implementa e executa os casos do seu domínio e entrega as evidências em
%% doc/tex/figuras/testes/ (ver DESENVOLVEDORES.md).
%% ===================================================================================
\chapter{Verificação, Validação e Testes}\label{cap:testes}

A \emph{verificação} responde se o sistema foi construído corretamente, isto é, conforme a especificação dos Capítulos~\ref{cap:requisitos} e~\ref{cap:arquitetura}; a \emph{validação} responde se foi construído o sistema correto, isto é, se atende às necessidades do usuário \cite{sommerville2011,pressman2021}. Este capítulo define como ambas serão realizadas e quais evidências serão coletadas.

\section{Objetivos e Escopo}\label{sec:testes-objetivos}

Os objetivos da verificação e validação são: (i) demonstrar que cada requisito essencial e importante é atendido (\req{RF01} a \req{RF26}); (ii) comprovar as regras de negócio invariantes (RN01 a RN17) no servidor; (iii) medir as qualidades declaradas como metas (\req{RNF02}, \req{RNF04}, \req{RNF05}); e (iv) verificar o comportamento do cache ISR em modo de produção. Ficam fora do escopo testes de carga em grande escala, pois o perfil de uso esperado é de um único administrador e tráfego de leitura moderado.

\section{Estratégia de Testes}\label{sec:testes-estrategia}

A estratégia segue a proporção clássica entre tipos de teste: muitos testes de unidade, menos testes de integração e poucos testes de ponta a ponta (\emph{E2E}).

\begin{table}[htb]
\caption{Tipos de teste, alvo e ferramentas adotadas}\label{tab:tipos-teste}
\small
\begin{tabularx}{\linewidth}{@{}L{2.6cm}YL{4.2cm}@{}}
\toprule
\textbf{Tipo} & \textbf{Alvo} & \textbf{Ferramentas adotadas (D-12)}\\
\midrule
Unidade & Funções puras: montagem da árvore, resolução de caminho, validadores, cálculo de contraste, \en{slug} & Vitest \cite{vitest}\\
\addlinespace
Integração & Funções de \arq{db.js} contra um SQLite temporário e isolado; Route Handlers com sessão simulada & Vitest com banco de teste (Seção~\ref{sec:testes-banco})\\
\addlinespace
Ponta a ponta & Fluxos completos no navegador sobre o build de produção (login, publicação, leitura, personalização) & Playwright \cite{playwright}\\
\addlinespace
Não funcionais & Desempenho, acessibilidade, segurança e ISR & Lighthouse \cite{lighthouse}, verificação de cabeçalhos, \texttt{npm audit}\\
\bottomrule
\end{tabularx}
\fonte{Elaborado pelos autores.}
\end{table}

\textbf{Ambiente de teste.} Cada suíte de integração usa um arquivo SQLite temporário, criado a partir das migrações, com \en{seed} controlado (Seção~\ref{sec:testes-banco}). Os testes de ISR e de desempenho usam \texttt{next build} seguido de \texttt{next start}, porque o comportamento de cache em \texttt{next dev} não representa a produção \cite{nextjsdocs}. Nenhum teste usa o banco de desenvolvimento nem segredos reais. O caso CT29 exige uma máquina com Docker e Compose, de preferência uma VPS limpa, para validar a instalação por imagem (\adr{ADR-011}).

\textbf{Ferramenta oficial.} O Vitest \cite{vitest} (versão 5, dependência de desenvolvimento) é a ferramenta oficial de testes de unidade e de integração, aprovada pelo Engenheiro Chefe (\decisao{D-12}, \adr{ADR-017}). O comando é \texttt{npm test} (\texttt{vitest run}) e a configuração está em \arq{vitest.config.mjs}, na raiz, com o mesmo alias \arq{@/} do \arq{jsconfig.json}. Os testes ficam em \arq{tests/}, por domínio: \arq{tests/setup/} (infraestrutura compartilhada), \arq{tests/dados/}, \arq{tests/core/} e \arq{tests/ui/}, que reutilizam \arq{tests/setup/}. Os testes de ponta a ponta usam o Playwright \cite{playwright}. O \arq{vitest.config.mjs} inclui um \emph{plugin} que transforma o JSX dos arquivos \texttt{.js} de \arq{src/} e de \arq{tests/}: a \adr{ADR-008} manda escrever JSX em \texttt{.js}, e o Vite só o transforma em \texttt{.jsx}; sem o \emph{plugin}, um teste não consegue importar um arquivo de rota nem um componente. Ele usa \texttt{transformWithOxc}, do próprio Vite 8 (dependência transitiva do Vitest, sem pacote novo), e não afeta o \texttt{.ts} do cliente gerado do Prisma; numa troca de versão do Vitest ou do Vite, esse \emph{import} precisa ser reavaliado.

\textbf{Responsabilidades.} Cada agente implementa os testes do seu domínio (Tabela~\ref{tab:casos-teste}); o Arquiteto verifica a cobertura dos requisitos e consolida os resultados no Capítulo~\ref{cap:resultados}.

\section{Testes de Integração com Banco Isolado}\label{sec:testes-banco}

Os testes que exercitam o banco seguem uma estratégia de isolamento, implementada pelo agente Dados em \arq{tests/setup/} (merge M-003), que todo teste de integração deve reutilizar:

\begin{enumerate}
  \item \textbf{Banco-modelo migrado.} Uma vez por rodada, o \texttt{globalSetup} cria um diretório temporário do sistema e aplica nele as migrações reais com \texttt{prisma migrate deploy}, o mesmo comando da produção. Assim, o esquema testado é exatamente o das migrações, inclusive o índice parcial da raiz (Seção~\ref{sec:arq-indice-parcial}), que o Prisma não modela.
  \item \textbf{Cópia por arquivo.} Antes de cada arquivo de teste, o \texttt{setupFiles} copia o banco-modelo para um arquivo próprio e define \texttt{DATABASE\_URL} com o caminho absoluto dessa cópia, antes de qualquer \texttt{import} de \arq{db.js} (que lê a variável ao ser carregado). Os arquivos de teste não interferem entre si.
  \item \textbf{Nunca o banco de desenvolvimento.} O \arq{dev.db} e o \arq{.env} não entram em jogo: a variável é sempre definida pelo \en{setup}, e uma salvaguarda aborta o teste se o banco estiver fora do diretório temporário. Ao final da rodada o diretório é removido.
  \item \textbf{Velocidade sem perder as constraints.} Nos bancos de teste, e só neles, usa-se \texttt{PRAGMA synchronous = OFF}; isso dispensa o \texttt{fsync} e não altera o comportamento das restrições verificadas.
  \item \textbf{Ausência de \en{drift}.} Os testes de migração conferem que \texttt{migrate status} e \texttt{migrate diff} não acusam divergência por causa do índice parcial escrito à mão.
\end{enumerate}

A primeira rodada, no merge M-003, tem 72 testes em cinco arquivos: restrições físicas do banco (8), \arq{content-tree} (14), conteúdo em \arq{db.js} (40), configuração (6) e migração com \en{seed} (4). Cobrem os casos CT11, CT12, CT13 e CT28 e três restrições da migração inicial: \en{slug} único na raiz, \en{slug} único entre irmãos e bloqueio da exclusão de usuário que possui conteúdo. Para validar a eficácia da suíte, o agente Dados introduziu mutações temporárias (remover o índice parcial, remover o índice composto, desativar a checagem de ciclo, desativar o filtro \texttt{PUBLISHED}, executar a reordenação sem transação) e confirmou que, em cada caso, os testes esperados falharam. Testes de Route Handlers devem definir \texttt{DATABASE\_URL} antes do primeiro \texttt{import} de \arq{db.js}, como faz \arq{tests/setup/per-file.js}.

O Core reutiliza essa infraestrutura em \arq{tests/core/} (merge M-004), com 269 testes em nove arquivos: validadores (93), utilitários HTTP e mapeamento de erros (33), revalidação (18), autenticação (16), comando \texttt{admin:create} (6), guarda de sessão (37), API de conteúdo (48), API de configuração (13) e Proxy (5). A diferença é que esses testes usam o \textbf{better-auth real}: o administrador é criado no banco de teste, o login passa por \texttt{auth.handler} e os Route Handlers recebem o \en{cookie} verdadeiro, de modo que a sessão é validada no banco, como em produção. O comando \texttt{admin:create} é executado de verdade com o \texttt{tsx}, com o \arq{.env} desligado. Apenas o \texttt{revalidatePath} é substituído por um duplo (\texttt{vi.mock("next/cache")}), pois exige o servidor do Next.js; o arquivo \arq{tests/core/env.js} define um segredo fictício. As recusas de escrita sem sessão (CT04) chamam os \en{handlers} diretamente, sem o Proxy, o que comprova que a barreira real está na API. Com as 72 do Dados, a suíte tinha 341 testes nesse merge. O Core também aplicou mutações temporárias (22, das quais 21 fizeram testes falharem; a restante era redundante), e o Arquiteto repetiu, na revisão, o \texttt{lint}, o \texttt{build} e a suíte completa.

A UI acrescenta, em \arq{tests/ui/} (merge M-005), 36 casos em um arquivo, \arq{theme.test.js}, sobre funções puras e sobre o \arq{globals.css}: \texttt{contrastRatio} (preto e branco = 21, cores iguais = 1, simetria, maiúsculas e minúsculas, valores conhecidos e recusa de entradas fora de \texttt{\#RRGGBB}), \texttt{contrastWarnings} (lista vazia; paleta padrão com avisos da primária e do destaque; limiar de 4,5:1 entre \texttt{\#767676} e \texttt{\#777777}; limiar de 3:1 do destaque; truncamento da razão) e \texttt{readableTextOn}, além da consistência dos tokens (a paleta padrão da RN08, o mapeamento \texttt{--color-*}, o valor de \arq{--vitrine-on-primary} igual ao de \texttt{readableTextOn(primária)}, o contraste dos neutros e a ausência de \texttt{prefers-color-scheme}). Cobrem a parte de unidade do CT24; a parte de API foi coberta pelo Core no merge M-006 (adiante). Com os 72 do Dados e os 269 do Core, a suíte tinha 377 testes em 15 arquivos nesse merge, reproduzidos pelo Arquiteto (\texttt{lint}, \texttt{build} e suíte completa). Não há testes automatizados de componentes, pois o ambiente de DOM não foi adotado: o comportamento da interface foi verificado no navegador pela UI (Seção~\ref{sec:impl-auth}), e o Playwright segue previsto para os casos de ponta a ponta.

O Core ampliou, no merge M-006, a API de configuração em \arq{tests/core/api-configuration.test.js}, de 13 para 23 testes (CT24, parte de API), de modo que \arq{tests/core/} passa a ter \textbf{279 testes em nove arquivos}. Os casos novos afirmam o campo, a razão truncada e o mínimo de cada aviso, e não só a presença da chave \texttt{warnings}: o destaque padrão (\texttt{\#f59e0b}) já gera aviso em qualquer \texttt{PUT} com a paleta padrão, e um teste genérico passaria pelo motivo errado. Cobrem a paleta padrão gravada (primária 3,67 e destaque 2,14, nessa ordem); o formato \texttt{\{ code, field, ratio, minimum, message \}}; a paleta sem falhas (\texttt{warnings} presente e vazio); a correção só da primária, que mantém o aviso do destaque; a secundária logo abaixo do limiar (4,47); a persistência dos avisos em um \texttt{PUT} posterior sem cores; o cálculo sobre as três cores \emph{lidas do banco}, e não sobre o corpo; a gravação e a revalidação normais com três avisos; e a falha no cálculo, que responde \texttt{500} depois de a configuração já estar gravada e o \en{layout} revalidado (resiliência; Seção~\ref{sec:arq-contraste}). O \texttt{GET} e a resposta \texttt{422} não trazem \texttt{warnings}.

Para observar e fazer falhar a função da UI sem alterá-la, o arquivo usa um \emph{espião transparente}: \texttt{vi.mock("@/lib/theme")} substitui \texttt{contrastWarnings} por um \texttt{vi.fn} que repassa a chamada à implementação real, o que permite conferir os argumentos e simular a falha com \texttt{mockImplementationOnce}, sem mudar o comportamento nos demais testes. O Core provou a eficácia dos testes com oito mutações temporárias da rota (calcular sobre o corpo; usar as cores padrão quando o corpo não traz cor; omitir \texttt{warnings} quando vazio; devolver sempre lista vazia; devolver \texttt{warnings} no \texttt{GET}; calcular antes de revalidar; responder 422 quando há aviso; avaliar só a primária), todas detectadas; o Arquiteto reproduziu três delas (calcular antes de revalidar, omitir \texttt{warnings} quando vazio e calcular sobre o corpo), também detectadas. Com os 72 do Dados, os 279 do Core e os 36 da UI, a suíte tinha 387 testes em 15 arquivos nesse merge, reproduzidos pelo Arquiteto (\texttt{lint}, \texttt{build} e suíte completa).

No merge M-007, o Core acrescentou 48 testes em dois arquivos novos de \arq{tests/core/}, de modo que o Core passa a ter \textbf{327 testes em 11 arquivos}. O \arq{site.test.js} (30) cobre o módulo \arq{src/app/_lib/site.js}: \texttt{buildThemeStyle} (as quatro variáveis; o texto sobre a primária igual a \texttt{readableTextOn(primária)}, comparado com a função e não com um literal, pois o texto escuro vai mudar para \texttt{\#000000}; doze entradas inválidas omitidas sem lançar; \texttt{null} dá objeto vazio), o nome e a descrição com os padrões e \texttt{getSiteConfiguration} (sem consulta quando \arq{NEXT_PHASE} vale \texttt{phase-production-build}; \texttt{null} quando a consulta lança). O \arq{layouts.test.js} (18) renderiza de verdade, com \texttt{renderToStaticMarkup} de \texttt{react-dom/server} e sobre o SQLite temporário, o \en{layout} raiz (\texttt{lang}; as quatro variáveis no \texttt{style} do \texttt{<html>} e nada no \texttt{<body>}; a paleta nova na renderização seguinte; ausência de \texttt{style} sem banco ou no \en{build}; cor inválida gravada direto no banco que não chega ao HTML), o \texttt{generateMetadata}, o \en{layout} do painel (a \en{prop} \texttt{siteName} e só ela; o nome no cabeçalho; a sessão validada \emph{antes} de ler a configuração) e a página 404 (texto em português, link para o início, nenhuma cor fixa do Tailwind, nenhuma menção a rascunho). Só se substituem os módulos que dependem do servidor do Next.js (\texttt{next/font/google}, a sessão do painel, a Server Function de sair e \texttt{usePathname}); o banco é o real, com um espião transparente em \texttt{getConfiguration}, como no M-006. O Core provou a eficácia dos testes com 21 mutações temporárias, todas detectadas; o Arquiteto reproduziu quatro (consultar no \en{build}, aceitar cor sem validação, ler a configuração antes da sessão e remover o \emph{plugin} de JSX), também detectadas.

Além dos testes, o CT23 foi verificado sem navegador por HTTP, sobre o \en{build} de produção com \texttt{next start}, pelo Core e, de forma independente, pelo Arquiteto: \arq{/login} e o painel trazem a paleta e o nome do banco; a 404 responde em português, com \texttt{noindex}; e, depois de um \texttt{PUT /api/configuration}, as páginas passam à nova paleta sem novo \en{build} (Seção~\ref{sec:arq-tema-estatico}). O \en{build} também passou com \texttt{DATABASE\_URL=file:/data/vitrine.db}, de diretório inexistente, e com um banco existente sem tabelas, sem criar pasta nem arquivo (o arquivo seguiu vazio). Com os 72 do Dados, os 327 do Core e os 36 da UI, a suíte tem \textbf{435 testes em 17 arquivos}, reproduzidos pelo Arquiteto (\texttt{lint}, \texttt{build} e suíte completa) na revisão do merge M-007.

\section{Casos de Teste}\label{sec:casos-teste}

A Tabela~\ref{tab:casos-teste} relaciona cada caso de teste aos requisitos e regras que ele verifica. Um caso é considerado \emph{aprovado} quando o resultado observado coincide com o esperado; um requisito essencial só é considerado atendido quando todos os seus casos de teste aprovam.

\begin{longtable}{@{}L{1.1cm}L{7.0cm}L{2.2cm}L{1.6cm}L{1.6cm}@{}}
\caption{Casos de teste e rastreabilidade}\label{tab:casos-teste}\\
\toprule
\textbf{ID} & \textbf{Caso de teste (ação e resultado esperado)} & \textbf{Verifica} & \textbf{Tipo} & \textbf{Agente}\\
\midrule
\endfirsthead
\toprule
\textbf{ID} & \textbf{Caso de teste (ação e resultado esperado)} & \textbf{Verifica} & \textbf{Tipo} & \textbf{Agente}\\
\midrule
\endhead
\bottomrule
\endfoot
CT01 & Login com credenciais válidas redireciona ao painel e cria a sessão & RF01 & E2E & Core\\[2pt]
CT02 & Login inválido retorna mensagem genérica, sem indicar o campo incorreto & RF01; RNF03 & API & Core\\[2pt]
CT03 & Acesso a \arq{/admin} sem sessão redireciona a \arq{/login} & RF01; RF21 & E2E & Core\\[2pt]
CT04 & Qualquer rota de escrita sem sessão retorna 401 & RF21; RN01 & API & Core\\[2pt]
CT05 & Encerrar sessão invalida a sessão no banco e remove o \en{cookie} & RF20 & API & Core\\[2pt]
CT06 & O cadastro público está desabilitado após o \en{bootstrap} & RF22 & API & Core\\[2pt]
CT07 & Criar documento válido retorna 201 e o item aparece na árvore & RF05; RF02 & API & Core\\[2pt]
CT08 & \en{Slug} duplicado entre irmãos retorna 409 (\arq{SLUG_CONFLICT}) & RN03 & API & Core\\[2pt]
CT09 & Pai que não é coleção retorna 422 (\arq{INVALID_PARENT}) & RN02 & API & Core\\[2pt]
CT10 & Mover um nó para dentro de um descendente é rejeitado (\arq{CYCLE_DETECTED}) & RN06 & API & Core\\[2pt]
CT11 & Reordenar irmãos persiste a nova ordem e a listagem a respeita & RF03; RN05 & Integração & Dados\\[2pt]
CT12 & Excluir uma coleção remove todos os descendentes & RF07; RN07 & Integração & Dados\\[2pt]
CT13 & Resolver caminho existente e publicado retorna o nó; inexistente, ou não publicado (RN15), retorna nulo (404) & RF14; RN04; RN15 & Unidade & Dados\\[2pt]
CT14 & Abrir uma coleção redireciona ao primeiro documento descendente & RN14 & E2E & UI\\[2pt]
CT15 & Fórmulas em linha e em bloco são renderizadas pelo KaTeX & RF09 & Unidade & UI\\[2pt]
CT16 & HTML bruto no Markdown é exibido como texto e não é executado & RN11; RNF03 & Unidade & UI\\[2pt]
CT17 & Tabelas, tachado e listas de tarefas (GFM) são renderizados & RF11 & Unidade & UI\\[2pt]
CT18 & Blocos de código recebem realce conforme a linguagem & RF10 & E2E & UI\\[2pt]
CT19 & A pré-visualização do editor é idêntica à página pública & RF08 & E2E & UI\\[2pt]
CT20 & O rascunho local sobrevive ao recarregamento da página & RF15 & E2E & UI\\[2pt]
CT21 & Publicar nova versão atualiza a página pública (revalidação da página) & RF16; RF17; RN10 & E2E & Core\\[2pt]
CT22 & Mudança estrutural (título, \en{slug}, pai, ordem) revalida o \en{layout} raiz & RF17; RN10 & E2E & Core\\[2pt]
CT23 & Alterar cores reflete-se nas páginas sem novo \en{build} & RF19 & E2E & UI\\[2pt]
CT24 & Cor fora do formato é recusada; contraste insuficiente gera aviso & RN09 & API/Unid. & Core, UI\\[2pt]
CT25 & Metadados de SEO presentes no HTML pré-renderizado & RF25; RNF07 & E2E & Core\\[2pt]
CT26 & A \en{sidebar} é navegável por teclado e destaca o item ativo & RF12; RNF05 & E2E & UI\\[2pt]
CT27 & O layout permanece utilizável em largura de 320 px & RNF06 & E2E & UI\\[2pt]
CT28 & A migração aplica-se em banco vazio e o \en{seed} é idempotente & RF23 & Integração & Dados\\[2pt]
CT29 & Uma instalação limpa em VPS com Docker, seguindo o Apêndice~\ref{ap:implantacao}, produz sistema funcional, e o banco persiste após recriar o contêiner & RNF01; D-04 & Manual & Core\\[2pt]
CT30 & O \en{backup} do banco (\texttt{VACUUM INTO}) pode ser restaurado com os dados íntegros & RNF10 & Manual & Dados\\[2pt]
CT31 & Um conteúdo \texttt{DRAFT}, e toda a sua subárvore, responde 404 por URL direta e não aparece na \en{sidebar}, na barra superior nem no mapa do site; o painel continua listando todos os conteúdos & RF26; RN15; RF25 & E2E & Core\\[2pt]
CT32 & A criação sem estado explícito resulta em \texttt{PUBLISHED}; alternar o estado revalida o \en{layout} raiz e o conteúdo passa a aparecer ou a desaparecer na visita seguinte, sem novo \en{build} & RN16; RF17; RF26 & E2E & Core\\[2pt]
CT33 & Link social com protocolo diferente de \texttt{http} e \texttt{https} é recusado (422); campo vazio é aceito e o ícone não é exibido; links válidos recebem \texttt{noopener} e \texttt{noreferrer} & RF24; RN17; RNF03 & API/Unid. & Core, UI\\[2pt]
\end{longtable}
\fonte{Elaborado pelos autores.}

\section{Testes Não Funcionais}\label{sec:testes-nf}

\subsection{Desempenho e ISR}

O desempenho será medido com o Lighthouse \cite{lighthouse} sobre o \en{build} de produção, em perfil móvel simulado e em desktop, registrando as métricas centrais de experiência \cite{webvitals}, em especial o \en{Largest Contentful Paint} (LCP). A meta é LCP $\leq$ 2,5 s no percentil 75 e pontuação de desempenho $\geq$ 90 (\req{RNF02}). O comportamento do cache será verificado por requisições repetidas à mesma página, observando o cabeçalho \arq{x-nextjs-cache}: a primeira requisição deve resultar em \texttt{MISS} ou \texttt{HIT} conforme a pré-renderização, as seguintes em \texttt{HIT}, e a primeira após uma mutação em \texttt{MISS} ou \texttt{REVALIDATED} (casos CT21 e CT22).

\vitrinefigura[tipo=teste, arquivo=teste-lighthouse-mobile, agente=UI, largura=0.85,
  descricao={Relatório do Lighthouse em perfil móvel para uma página de documento no build de produção, com as quatro pontuações e as métricas LCP, CLS e TBT.}]
  {Relatório do Lighthouse (perfil móvel)}{fig:teste-lh-mobile}

\vitrinefigura[tipo=teste, arquivo=teste-isr-cabecalhos, agente=Core, largura=0.85,
  descricao={Saída de terminal com requisições repetidas e o cabeçalho x-nextjs-cache (MISS, HIT, HIT, e após mutação REVALIDATED ou MISS).}]
  {Evidência do comportamento do cache ISR}{fig:teste-isr}

\subsection{Segurança}

A verificação de segurança consiste em: executar os casos CT02, CT04, CT06, CT16, CT31 e CT33; conferir os cabeçalhos de segurança da resposta; executar \texttt{npm audit} e registrar as vulnerabilidades encontradas (os quatro alertas altos já avaliados e aceitos constam na Seção~\ref{sec:arq-npm-audit} e na \adr{ADR-019}; qualquer alerta novo exige nova avaliação); e percorrer a lista da Tabela~\ref{tab:ameacas}, confirmando cada controle. Como referência de itens de verificação utiliza-se o ASVS do OWASP \cite{owaspasvs}.

\subsection{Acessibilidade}

A acessibilidade será avaliada segundo o WCAG 2.1 nível AA \cite{wcag21}, combinando verificação automatizada (contraste, rótulos, estrutura de títulos) e verificação manual (navegação apenas por teclado, foco visível, leitura por leitor de tela em um fluxo representativo).

\vitrinefigura[tipo=teste, arquivo=teste-acessibilidade, agente=UI, largura=0.85,
  descricao={Resultado da verificação automatizada de acessibilidade (por exemplo, categoria Acessibilidade do Lighthouse ou ferramenta equivalente) para as páginas pública e do painel.}]
  {Resultado da verificação automatizada de acessibilidade}{fig:teste-a11y}

\subsection{Usabilidade}

A usabilidade será validada com usuários do perfil do Administrador (docentes e pesquisadores), em sessões individuais com tarefas roteirizadas e o questionário SUS \cite{brooke1996sus}. Estudos de usabilidade clássicos indicam que amostras pequenas já revelam a maior parte dos problemas de uma interface \cite{nielsen1993usability}; planeja-se um grupo de cinco a oito participantes. As tarefas previstas são:

\begin{enumerate}[label=\textbf{T\arabic*.}, leftmargin=1.2cm]
  \item entrar no painel e criar uma coleção chamada ``Cálculo'';
  \item criar, dentro dela, um documento com um título, um parágrafo e uma fórmula;
  \item conferir que o documento está publicado e encontrá-lo na página pública; colocá-lo em rascunho e confirmar que deixa de aparecer;
  \item mover o documento para outra coleção e reordená-lo;
  \item alterar a cor primária do site e confirmar o efeito na página pública;
  \item preencher o link do GitHub nas configurações e localizar o ícone correspondente no rodapé.
\end{enumerate}

A meta é pontuação SUS média igual ou superior a 68 (\req{RNF04}) e taxa de conclusão de pelo menos 80\% por tarefa.


\section{Critérios de Saída}\label{sec:testes-saida}

A versão é considerada pronta para a entrega (\en{release}) quando: (i) todos os casos de teste dos requisitos \emph{Essenciais} aprovam; (ii) os casos dos requisitos \emph{Importantes} aprovam ou têm justificativa de exceção registrada; (iii) \texttt{npm run lint}, \texttt{npm run build} e \texttt{npm test} passam sem erros; (iv) não há vulnerabilidades críticas pendentes no relatório de auditoria de dependências, e as vulnerabilidades altas se limitam às aceitas e documentadas na \adr{ADR-019}; (v) as metas de \req{RNF02}, \req{RNF04} e \req{RNF05} foram medidas e registradas, ainda que não atingidas, com análise no Capítulo~\ref{cap:resultados}.

\vitrinefigura[tipo=teste, arquivo=teste-suite-automatizada, agente=Core, largura=0.85,
  descricao={Saída do executor de testes (unidade, integração e E2E) mostrando o total de casos aprovados, falhos e ignorados e o tempo de execução.}]
  {Execução da suíte de testes automatizados}{fig:teste-suite}

\section{Síntese do Capítulo}

O capítulo definiu a estratégia de verificação e validação, trinta e três casos de teste rastreados a requisitos e regras, os procedimentos não funcionais (desempenho e ISR, segurança, acessibilidade e usabilidade) e os critérios de saída. A execução e a análise dos resultados são apresentadas no capítulo seguinte.
